\documentclass[10pt,a4paper]{amsart}
\usepackage{amsmath,amssymb,amsthm,graphicx}
\usepackage{verbatim}
\newenvironment{reference}{\comment}{\endcomment}
%\newenvironment{reference}{}{}
\newenvironment{slogan}{\comment}{\endcomment}
\newenvironment{history}{\comment}{\endcomment}

% For commutative diagrams we use Xy-pic
\usepackage[all]{xy}

% We use 2cell for 2-commutative diagrams.
\xyoption{2cell}
% Theorem environments.
%
\theoremstyle{plain}
\newtheorem{theorem}[subsection]{Theorem}
\newtheorem{proposition}[subsection]{Proposition}
\newtheorem{lemma}[subsection]{Lemma}

\theoremstyle{definition}
\newtheorem{definition}[subsection]{Definition}
\newtheorem{example}[subsection]{Example}
\newtheorem{exercise}[subsection]{Exercise}
\newtheorem{situation}[subsection]{Situation}

\theoremstyle{remark}
\newtheorem{remark}[subsection]{Remark}
\newtheorem{remarks}[subsection]{Remarks}

\numberwithin{equation}{subsection}

% Macros
%
\def\lim{\mathop{\mathrm{lim}}\nolimits}
\def\colim{\mathop{\mathrm{colim}}\nolimits}
\def\Spec{\mathop{\mathrm{Spec}}}
\def\Hom{\mathop{\mathrm{Hom}}\nolimits}
\def\Ext{\mathop{\mathrm{Ext}}\nolimits}
\def\SheafHom{\mathop{\mathcal{H}\!\mathit{om}}\nolimits}
\def\SheafExt{\mathop{\mathcal{E}\!\mathit{xt}}\nolimits}
\def\Sch{\mathit{Sch}}
\def\Mor{\mathop{\mathrm{Mor}}\nolimits}
\def\Ob{\mathop{\mathrm{Ob}}\nolimits}
\def\Sh{\mathop{\mathit{Sh}}\nolimits}
\def\NL{\mathop{N\!L}\nolimits}
\def\CH{\mathop{\mathrm{CH}}\nolimits}
\def\proetale{{pro\text{-}\acute{e}tale}}
\def\etale{{\acute{e}tale}}
\def\QCoh{\mathit{QCoh}}
\def\Ker{\mathop{\mathrm{Ker}}}
\def\Im{\mathop{\mathrm{Im}}}
\def\Coker{\mathop{\mathrm{Coker}}}
\def\Coim{\mathop{\mathrm{Coim}}}

% Boxtimes
%
\DeclareMathSymbol{\boxtimes}{\mathbin}{AMSa}{"02}

%
% Macros for moduli stacks/spaces
%
\def\QCohstack{\mathcal{QC}\!\mathit{oh}}
\def\Cohstack{\mathcal{C}\!\mathit{oh}}
\def\Spacesstack{\mathcal{S}\!\mathit{paces}}
\def\Quotfunctor{\mathrm{Quot}}
\def\Hilbfunctor{\mathrm{Hilb}}
\def\Curvesstack{\mathcal{C}\!\mathit{urves}}
\def\Polarizedstack{\mathcal{P}\!\mathit{olarized}}
\def\Complexesstack{\mathcal{C}\!\mathit{omplexes}}
% \Pic is the operator that assigns to X its picard group, usage \Pic(X)
% \Picardstack_{X/B} denotes the Picard stack of X over B
% \Picardfunctor_{X/B} denotes the Picard functor of X over B
\def\Pic{\mathop{\mathrm{Pic}}\nolimits}
\def\Picardstack{\mathcal{P}\!\mathit{ic}}
\def\Picardfunctor{\mathrm{Pic}}
\def\Deformationcategory{\mathcal{D}\!\mathit{ef}}
\usepackage{cleveref}
\crefname{originalreference}{original reference}{original references}
\Crefname{originalreference}{Original reference}{Original references}
\makeatletter
\newlabel{sites-lemma-recover}{{00VD}{source}{}{sourcepointer-1}{}}
\newlabel{sites-lemma-recover@cref}{{[originalreference][0][]00VD}{[1][0][]source}{}{}{}}
\makeatother
\title{Adjunction formulas for the pushforward and pullback functors along a morphism of sites}
\author{The Stacks Project Authors}
\date{stacks-project:sites-modules.tex:local-snapshot-20261012; GNU FDL 1.2 or later}
\begin{document}
\maketitle
\noindent\textit{Editorial note.} All delivered mathematical text is the author's own, copied verbatim from the pinned Stacks Project chapter file sites-modules.tex; no statement or proof step was rewritten, reordered or removed.\par\medskip
\noindent\textit{Editorial note.} The counted span cites one external item (a section of the same chapter family), recorded in assembly.external\_references with its canonical tag as a navigation pointer.\par\medskip
\noindent\textit{Allowed standard prerequisites (editorial).}\begin{itemize}
\item Definitions of g\_!, g\textasciicircum{}\{-1\}, u\textasciicircum{}p and g\_\{p!\} for a morphism of sites. Stacks Project, Modules on Sites, section on morphisms of sites (definitions of the four functors). Supplies the functors whose adjunction is being established.
\item Universal property of sheafification. Stacks Project, Sites, section on sheaves on a site. Converts the presheaf Hom-set into a sheaf Hom-set, giving the second formula.
\end{itemize}\medskip
\ifdefined\hypertarget\hypertarget{sourcepointer-1}{}\fi
\noindent\textit{Source reference.} Original reference 00VD: sources/stacks-chow/tags.txt, source lines 1143--1143. Cross-chapter standard prerequisite cited by the author inside the counted span; it is used as a pointer and its own proof is not part of this item.\par\medskip
\begin{lemma}
\label{lemma-g-shriek-adjoint}
The functor $g_{p!}$ is a left adjoint to the functor $u^p$.
The functor $g_!$ is a left adjoint to the functor $g^{-1}$.
In other words the formulas
\begin{align*}
\Mor_{\textit{PAb}(\mathcal{C})}(\mathcal{F}, u^p\mathcal{G})
& =
\Mor_{\textit{PAb}(\mathcal{D})}(g_{p!}\mathcal{F}, \mathcal{G}), \\
\Mor_{\textit{Ab}(\mathcal{C})}(\mathcal{F}, g^{-1}\mathcal{G})
& =
\Mor_{\textit{Ab}(\mathcal{D})}(g_!\mathcal{F}, \mathcal{G})
\end{align*}
hold bifunctorially in $\mathcal{F}$ and $\mathcal{G}$.
\end{lemma}

\begin{proof}
The second formula follows formally from the first, since if
$\mathcal{F}$ and $\mathcal{G}$ are abelian sheaves then
\begin{align*}
\Mor_{\textit{Ab}(\mathcal{C})}(\mathcal{F}, g^{-1}\mathcal{G})
& =
\Mor_{\textit{PAb}(\mathcal{D})}(g_{p!}\mathcal{F}, \mathcal{G}) \\
& =
\Mor_{\textit{Ab}(\mathcal{D})}(g_!\mathcal{F}, \mathcal{G})
\end{align*}
by the universal property of sheafification.

\medskip\noindent
To prove the first formula, let $\mathcal{F}$, $\mathcal{G}$ be abelian
presheaves. To prove the lemma we will construct maps from the group on the
left to the group on the right and omit the verification that these are
mutually inverse.

\medskip\noindent
Note that there is a canonical map of abelian presheaves
$\mathcal{F} \to u^pg_{p!}\mathcal{F}$ which on sections over $U$ is the
natural map
$\mathcal{F}(U) \to \colim_{u(U) \to u(U')} \mathcal{F}(U')$, see
Sites, Lemma \ref{sites-lemma-recover}.
Given a map $\alpha : g_{p!}\mathcal{F} \to \mathcal{G}$
we get $u^p\alpha : u^pg_{p!}\mathcal{F} \to u^p\mathcal{G}$.
which we can precompose by the map $\mathcal{F} \to u^pg_{p!}\mathcal{F}$.

\medskip\noindent
Note that there is a canonical map of abelian presheaves
$g_{p!}u^p\mathcal{G} \to \mathcal{G}$ which on sections over
$V$ is the natural map
$\colim_{V \to u(U)} \mathcal{G}(u(U)) \to \mathcal{G}(V)$.
It maps a section $s \in u(U)$ in the summand corresponding to
$t : V \to u(U)$ to $t^*s \in \mathcal{G}(V)$.
Hence, given a map $\beta : \mathcal{F} \to u^p\mathcal{G}$
we get a map $g_{p!}\beta : g_{p!}\mathcal{F} \to g_{p!}u^p\mathcal{G}$
which we can postcompose with the map $g_{p!}u^p\mathcal{G} \to \mathcal{G}$
above.
\end{proof}

\end{document}
